\begin{abstract}
Large-scale unsupervised language models (LMs) are adept at acquiring general world knowledge and certain reasoning abilities; however, it remains challenging to exercise fine-grained control over their outputs because their training occurs without supervision. Current approaches to increase such controllability typically involve collecting human judgments about the relative merits of model outputs and then fine-tuning the unsupervised LM to align with these labeled preferences, frequently via reinforcement learning from human feedback (RLHF). Despite its effectiveness, RLHF is a multi-stage and sometimes unstable approach: it involves initially training a reward model to capture human preferences and subsequently employing reinforcement learning to adapt the unsupervised LM so as to maximize the inferred reward, while attempting to prevent divergence from the pretrained LM. In this work, we propose a novel reparameterization of the RLHF reward model that makes it possible to directly obtain the optimal policy in closed form, thereby reducing the conventional RLHF optimization process to a straightforward classification task. This leads to our method, named \textit{Direct Preference Optimization} (DPO), which is stable, efficient, and easy to compute—removing the need for sampling from the LM during fine-tuning or extensive tuning of hyperparameters. Through our empirical studies, we demonstrate that DPO is capable of fine-tuning LMs to better conform to human preferences, achieving results that are on par with or superior to established techniques. Specifically, DPO-based fine-tuning surpasses RLHF with PPO in sentiment control and matches or improves upon existing systems in summarization and single-turn dialogue tasks, all while substantially simplifying the implementation and training process.
\end{abstract}

\section{Introduction}
Extensive unsupervised training of language models (LMs) on massive corpora has resulted in emergent abilities that exceed expectations~\citep{chowdhery2022palm, brown2020language, touvron2023llama, bubeck2023sparks}. Nevertheless, because the training data reflect diverse human intentions, preferences, and expertise, not every trait present in the dataset is appropriate for models to emulate. For instance, while an AI coding assistant should be capable of \textit{recognizing} widespread programming errors to provide corrective suggestions, it is preferable that it prioritize emulating the subset of exemplary programming evident in the data. Likewise, though a model ought to be \textit{cognizant} of widely held misconceptions—such as those accepted by 50\% of people—it is important that it not endorse those false beliefs in half of the relevant outputs. Thus, the process of \emph{determining which behaviors and answers are preferred} among the model’s extensive repertoire of \textit{knowledge and competencies} is fundamental for developing AI systems that are reliable, effective, and governable \citep{ouyang2022training}. Existing strategies typically employ reinforcement learning (RL) to guide language models towards alignment with human preferences. In this paper, we demonstrate that the commonly used RL objective for preference learning can actually be exactly matched by a straightforward binary cross-entropy loss, thereby streamlining the learning process.

Generally, current approaches integrate the desired attributes into language models by utilizing tailored collections of human preference data, reflecting behaviors valued for safety and utility. Preference learning is carried out as a separate phase, subsequent to an initial unsupervised pre-training on a vast text corpus. Although supervised fine-tuning using human-generated high-quality response examples is the most direct approach to preference alignment, reinforcement learning from human or AI feedback (RLHF/RLAIF; \citep{christiano2017deep,bai2022constitutional}) has shown greater empirical success. In RLHF frameworks, a reward model is trained on human preference datasets, and RL is then used to update the language model policy to prefer outputs that yield higher rewards, all while minimizing deviation from the pretrained model. While RLHF methodologies yield language models with advanced capabilities in conversation and coding, they introduce substantial additional complexity relative to standard supervised methods; these workflows typically require the training of several LMs and demand costly policy sampling during the training loop.

This work presents a method for optimizing language models directly in accordance with human preferences, foregoing the need for explicit reward models or reinforcement learning techniques. We introduce \textit{Direct Preference Optimization (DPO)}, an approach that, while targeting the same objective as standard RLHF algorithms—namely, maximizing reward subject to a KL-divergence penalty—is both straightforward to implement and easy to train. DPO intuitively operates by raising the relative log-likelihood of preferred outputs versus non-preferred ones, incorporating a per-example, dynamically determined importance weighting. This mechanism safeguards against model collapse, which we observe when naively optimizing a probability ratio. Like prior algorithms, DPO leverages a theoretical preference model (such as the Bradley-Terry model; \cite{bradley1952rankanalysis}) to quantify the alignment between a reward function and observed preference judgments. However, whereas conventional methods construct a preference loss via the preference model to supervise reward model training before subsequently training a policy on this learned reward, DPO applies a variable substitution to directly express the preference loss in terms of the policy itself. Provided a dataset reflecting human choices over alternative outputs, DPO can thus optimize the policy using a binary cross-entropy loss, effectively recovering the optimal policy under an implicit reward structure inferred from the observed preferences.

The core contribution of this study is the Direct Preference Optimization (DPO) algorithm—a reinforcement learning-free framework for preference-based language model training. Empirical evaluations reveal that DPO matches or exceeds the performance of established methods, such as RLHF implemented with PPO, across preference-driven tasks including sentiment control, summarization, and dialogue generation, tested on models up to 6B parameters in size.

As self-supervised language models have grown in size, their ability to address tasks with zero-shot \citep{radford2019language} or few-shot prompting \citep{gpt3,megatron,chowdhery2022palm} has notably improved. Nevertheless, their effectiveness on specific downstream applications and adherence to user intentions is considerably enhanced through the process of fine-tuning on corpora composed of instructions paired with human-generated responses \citep{mishra-etal-2022-cross,sanh2022multitask,chung2022scaling,thoppilan2022lamda}. This technique, often referred to as `instruction-tuning', not only broadens the scope of instructions to which LLMs can effectively generalize but also typically results in greater practicality and utility \citep{chung2022scaling}. In spite of its advantages, it is often easier to collect \textit{relative} human assessments of response quality than expert-level demonstration data; consequently, more recent approaches have involved fine-tuning LLMs using collections of human preference data, leading to advancements in translation \citep{kreutzer-etal-2018-reliability}, summarization \citep{stiennon2022learning,ziegler2020finetuning}, narrative generation \citep{ziegler2020finetuning}, and following instructions \citep{ouyang2022training,ramamurthy2023is}. These strategies typically start by learning a neural reward function that fits the human preference data, often employing a preference framework such as the Bradley-Terry model \citep{bradley1952rankanalysis}. The language model is then further trained to maximize this reward via reinforcement learning algorithms such as REINFORCE \citep{williams1992reinforce}, proximal policy optimization (PPO; \cite{schulman2017proximal}), or adaptations thereof \citep{ramamurthy2023is}. Another related stream of research uses LLMs fine-tuned through human feedback on instructions to generate extra synthetic preference examples aimed at properties like safety or harmlessness \citep{bai2022constitutional}, utilizing limited human supervision through text-based criteria for annotation. These developments mark a confluence between research on reinforcement learning-based language model training for various targets~\citep{Ranzato2015SequenceLT,paulus2018a,wu2018learning} and general strategies for preference-based learning from humans \citep{christiano2017deep,kupcsik2018learning}. Despite the attractiveness of leveraging relative human preferences, the practicalities of applying reinforcement learning for fine-tuning large language models present significant challenges; the present work introduces a theoretically principled method for optimizing models based on relative preferences without reliance on reinforcement learning.

Beyond linguistic applications, learning policies based on preferences has been explored in both bandit and reinforcement learning paradigms, with diverse methodologies introduced in the literature. One notable line of work is contextual dueling bandit (CDB; \cite{yue2012karmed,dudik2015contextual}), which extends contextual bandit learning to scenarios where preferences or rankings over actions are available in lieu of explicit rewards. Theoretical analyses in this context, where absolute rewards are lacking, replace the standard definition of an optimal policy with the concept of a \textit{von Neumann winner}. Such a policy guarantees an expected win rate of at least 50\% against \textit{every} other policy \citep{dudik2015contextual}. Unlike CDBs, which acquire preference information online, human preference learning commonly relies on pre-collected, static datasets consisting of annotated preference pairs \citep{yan2022human}. In a similar vein, \textit{preference-based RL} (PbRL) frameworks derive supervision from binary preferences generated by an \textit{unknown} scoring mechanism instead of reward signals \citep{BusaFekete2014,ruiz2023dueling}. Several techniques have been developed for PbRL, many of which support reusing off-policy preference data; nevertheless, these generally entail a two-stage approach where the unknown scoring function (i.e., the reward model) is first inferred, followed by policy optimization with respect to this estimated function \citep{jain2013learning,BusaFekete2014,christiano2017deep,sadigh2017active,kupcsik2018learning}. By contrast, we introduce a unified, single-stage strategy for policy learning, which aims to directly optimize the policy in accordance with observed preferences.

\section{Preliminaries}\label{section:prelims}

We now summarize the RLHF procedure as outlined in \citeauthor{ziegler2020finetuning}, as well as subsequent elaborations \citep{stiennon2022learning, bai2022training, ouyang2022training}. This pipeline generally consists of three key stages: (1) supervised fine-tuning (SFT); (2) preference collection and reward modeling; and (3) reinforcement learning-based policy optimization.

\textbf{SFT}: The RLHF process typically starts with the supervised fine-tuning of a large language model that has been previously pre-trained. This step leverages curated, high-quality labeled data aligned with specific downstream objectives (such as dialogue generation, summarization, and so forth), resulting in a model denoted by $\pi^\text{SFT}$.

\textbf{Reward Modelling Phase}: In the subsequent phase, the SFT model is provided prompts $x$ and outputs two candidate completions $(y_1, y_2)\sim \pi^\text{SFT}(y \mid x)$. These output pairs are presented to human annotators, who indicate their preference between the completions—this is noted as $y_w\succ y_l \mid x$, with $y_w$ and $y_l$ identifying the preferred and non-preferred responses, respectively. These expressed preferences are presumed to arise from some underlying, unobserved reward function $r^*(y, x)$. Several frameworks have been proposed to mathematically capture such preferences, among which the Bradley-Terry (BT) \cite{bradley1952rankanalysis} model is prevalent (though more general models such as the Plackett-Luce ranking model \citep{plackett1975analysis, luce2012individual} can also be accommodated if multiple ranked choices are provided). According to the BT model, the probability distribution $p^*$ representing human preference is expressed as:

\begin{equation}\label{eq:bradley-terry}
    p^*(y_1\succ y_2 \mid x)=\frac{\exp\left(r^*(x, y_1)\right)}{\exp\left(r^*(x, y_1)\right) + \exp\left(r^*(x, y_2)\right)}.
\end{equation}

Given a fixed dataset of preference comparisons $\mathcal{D}=\bigl\{x^{(i)}, y_w^{(i)}, y_l^{(i)}\bigr\}_{i=1}^N$ drawn according to $p^*$, we introduce a parameterized reward function $r_{\phi}(x, y)$ and fit its parameters via maximum likelihood estimation. When viewed as a binary classification problem, the negative log-likelihood objective becomes:

\begin{equation}\label{eq:reward_model}
    \mathcal{L}_R(r_{\phi}, \mathcal{D}) = -\mathbb{E}_{(x, y_w, y_l)\sim \mathcal{D}}\bigl[\log \sigma(r_{\phi}(x, y_w)- r_{\phi}(x, y_l))\bigr]
\end{equation}

with $\sigma$ denoting the logistic sigmoid. For language models, it is common to initialize $r_{\phi}(x, y)$ from the SFT policy $\pi^\text{SFT}(y \mid x)$, appending a linear transformation atop the final transformer layer to produce the scalar reward output~\cite{ziegler2020finetuning}. To minimize variance in the reward outputs, prior studies normalize rewards, enforcing $\mathbb{E}_{x,y\sim \mathcal{D}}\left[r_\phi(x, y)\right] = 0$ across $x$.

\textbf{RL Fine-Tuning Phase}: In the subsequent RL fine-tuning step, the acquired reward model is utilized to generate learning signals for the language model. Following previous frameworks~\citep{jaques2017sequence, jaques2020human}, the optimization target is defined as

\begin{equation}\label{eq:RL}
\max_{\pi_{\theta}}  \mathbb{E}_{x\sim \mathcal{D}, y\sim \pi_{\theta}(y \mid x)}\bigl[r_{\phi}(x, y)\bigr] - \beta\mathbb{D}_{\textrm{KL}}\bigl[\pi_{\theta}(y\mid x)\mid \mid \pi_\text{ref}(y\mid x)\bigr],
\end{equation}

where $\beta$ controls the extent to which the learned policy is regularized against the base reference model $\pi_\text{ref}$—often set as the original SFT policy $\pi^\text{SFT}$. 
In implementation, $\pi_\theta$ is generally initialized with $\pi^\text{SFT}$ as well. This KL penalty ensures that policy updates remain within regions where the reward model remains reliable and sustains a degree of diversity in generated outputs, thereby averting collapse to a narrow set of responses. Owing to the discreteness of language generation, this objective lacks differentiability and is typically addressed using reinforcement learning algorithms. The widely used strategy~\cite{ziegler2020finetuning, stiennon2022learning, bai2022training, ouyang2022training} introduces the reward function ${r(x, y) = r_{\phi}(x, y) -\beta (\log \pi_{\theta}(y\mid x) - \log \pi_\text{ref}(y\mid x))}$, and optimizes it with PPO~\cite{schulman2017proximal}. 

\section{Direct Preference Optimization}\label{sec:DPO}

Seeking to circumvent the computational overhead and practical complexity of applying reinforcement learning in large-scale language model fine-tuning, we aim to devise a direct, preference-based method for policy improvement. Unlike standard RLHF procedures that first estimate a reward function and subsequently perform RL optimization, our methodology hinges on a distinct parameterization of the reward model, which allows the associated optimal policy to be explicitly characterized in closed form—eliminating the need for a full RL loop. As elaborated in the following sections, our main insight involves exploiting an analytic mapping that links reward models to their optimal policies, thus enabling conversion of loss functions over rewards into equivalent policy-based losses. This reparameterization allows us to optimize policies directly under human-preference models such as the Bradley-Terry formulation, without explicitly constructing a standalone reward function. Here, the policy network effectively subsumes the roles of both the language generator and the (

\textbf{Reformulating the DPO objective.} Our starting point aligns with established literature: we consider the reinforcement learning (RL) objective presented in Eq.~\ref{eq:RL}, where the reward function $r$ is arbitrary. Drawing on previous research~\citep{peters2007reinforcement, peng2019advantage, korbak2022reinforcement, go2023aligning}, it follows directly that the optimal policy for maximizing expected reward under a Kullback-Leibler (KL) divergence constraint—specifically, Eq.~\ref{eq:RL}—is given by:

\begin{equation}\label{eq:op_policy}
    \pi_r(y\mid x) = \frac{1}{Z(x)}\pi_\text{ref}(y\mid x)\exp\left(\frac{1}{\beta}r(x, y)\right),
\end{equation}

with $Z(x) =\sum_{y}\pi_\text{ref}(y\mid x)\exp\left(\frac{1}{\beta}r(x, y)\right)$ representing the normalizing constant, or partition function. A full derivation is included in Appendix \ref{app:derivation1}. Notably, even substituting the reward estimator $r_{\phi}$ for the true reward $r^*$ leads to substantial computational difficulty when evaluating $Z(x)$~\citep{korbak2022reinforcement, go2023aligning}, limiting practical use of this formula. To circumvent this challenge, Eq.~\ref{eq:op_policy} can be algebraically rearranged to rewrite the reward $r(x, y)$ in terms of the optimal policy $\pi_r$, the reference policy $\pi_\text{ref}$, and $Z(\cdot)$. By applying the logarithm to both sides of Eq.~\ref{eq:op_policy}, we derive:

\begin{equation}\label{eq:main_eq}
    r(x,y) =\beta \log \frac{\pi_r(y\mid x)}{\pi_\text{ref}(y\mid x)} + \beta \log Z(x).
\end{equation}

This change of variables is applicable to both the ground-truth reward $r^*$ and its corresponding optimal policy $\pi^*$. A crucial simplification arises from the structure of the Bradley-Terry model: it depends solely on reward differences for two completions, namely, ${p^*(y_1 \succ y_2 \mid x) = \sigma(r^*(x, y_1) - r^*(x, y_2))}$. By substituting the above form (Eq.~\ref{eq:main_eq}) of $r^*(x, y)$ into the preference model (Eq.~\ref{eq:bradley-terry}), the partition function $Z(x)$ cancels out, enabling us to express the human preference likelihood in terms only of $\pi^*$ and $\pi_\text{ref}$. As a result, the optimal RLHF policy $\pi^*$ under the Bradley-Terry setup adheres to:

\begin{equation}\label{eq:objective}
    p^*(y_1\succ y_2 \mid x)=\frac{1}{1 + \exp\left(\beta \log \frac{\pi^*(y_2\mid x)}{\pi_\text{ref}(y_2\mid x)} - \beta \log \frac{\pi^*(y_1\mid x)}{\pi_\text{ref}(y_1\mid x)}\right)}
\end{equation}

The step-by-step details are available in Appendix~\ref{app:derivation2}. Although Eq.~\ref{eq:objective} is framed in terms of the Bradley-Terry model, similar derivations apply to more general settings such as the Plackett-Luce family~\citep{plackett1975analysis, luce2012individual}, as illustrated in Appendix~\ref{app:plackett_luce_models}.

With this formulation, preference probabilities from human annotation can now be re-expressed in terms of the optimal policy rather than relying on an explicit reward function. This insight leads us to a maximum likelihood objective defined for a parameterized policy $\pi_\theta$. Mirroring the logic of reward modeling (see Eq.~\ref{eq:reward_model}), our policy-based loss for DPO becomes:

\begin{equation}\label{eq:optimum_model}
    \mathcal{L}_\text{DPO}(\pi_{\theta}; \pi_\text{ref}) = -\mathbb{E}_{(x, y_w, y_l)\sim \mathcal

In this approach, we learn an implicit reward function via an alternative parametrization, where the optimal policy corresponds directly to $\pi_\theta$. Furthermore, as our methodology aligns with fitting a reparameterized Bradley-Terry model, it inherits theoretical guarantees, such as consistency provided that the preference data distribution satisfies certain conditions \cite{bong2022generalized}. We elaborate on the theoretical aspects of DPO and its relationship to related literature in Section~\ref{sec:theory}.

\textbf{Interpreting the DPO update.} To better grasp the internal workings of DPO, we examine the gradient of the loss $\mathcal{L}_\text{DPO}$. Specifically, the gradient with respect to $\theta$ takes the form:

\begin{multline*}\label{eq:gradient}
    \nabla_\theta \mathcal{L}_\text{DPO}(\pi_\theta;\pi_\text{ref}) = \\ -\beta\mathbb{E}_{(x, y_w, y_l) \sim \mathcal{D}} \bigg[\underbrace{\sigma(\hat{r}_\theta(x, y_l) - \hat{r}_\theta (x, y_w))}_\text{greater weight when the reward estimate errs}\bigg[\underbrace{\nabla_\theta\log \pi(y_w \mid x)}_\text{increase probability of $y_w$} - \underbrace{\nabla_\theta\log\pi(y_l \mid x)}_\text{reduce probability of $y_l$}\bigg]\bigg],
\end{multline*}

where $\hat{r}_\theta(x, y) = \beta \log \frac{\pi_\theta(y \mid x)}{\pi_\text{ref}(y \mid x)}$ serves as the implicit reward function determined by the language model $\pi_\theta$ relative to the reference model $\pi_\text{ref}$ (see Section~\ref{sec:theory} for further discussion). Conceptually, the loss gradient for $\mathcal{L}_\text{DPO}$ acts to boost the likelihood of favored responses $y_w$, while diminishing the likelihood of less preferred responses $y_l$. Crucially, the influence of each example depends on the degree to which the implicit reward model $\hat{r}_\theta$ overrates the dispreferred option, with scaling by $\beta$ to capture both the error in the ranking and the KL regularization strength. Our empirical results highlight the necessity of this weighting; omitting the weighting factor can result in model collapse, as evidenced in Appendix Table~\ref{tab:unlikelihood_generations}.

\textbf{Overview of DPO.} The typical DPO workflow includes: 1) Drawing response samples $y_1, y_2 \sim \pi_\text{ref}(\cdot \mid x)$ for each input prompt $x$, then annotating them based on human judgments to assemble an offline preference dataset $\mathcal{D} = \{x^{(i)}, y_w^{(i)}, y_l^{(i)}\}_{i=1}^N$; and 2) training the language model $\pi_\theta$ by minimizing $\mathcal{L}_\text{DPO}$, given $\pi_\text{ref}$, $\mathcal{D}$, and the specified value of $\beta$. 
In real-world applications, practitioners typically opt to leverage available public preference datasets instead of producing new samples and collecting human feedback. Because these datasets are generated using $\pi^\text{SFT}$, whenever possible we set $\pi_\text{ref} = \pi^\text{SFT}$ for initialization. In cases where $\pi^\text{SFT}$ is missing, $\pi_\text{ref}$ is initialized by fitting it to the distribution of preferred outputs ${(x, y_w)}$; concretely, this involves solving ${\pi_\text{ref} = \argmax_{\pi}\mathbb{E}_{x, y_w \sim \mathcal{D}}\left[\log \pi(y_w \mid x)\right]}$. Such initialization reduces discrepancies between the actual (but intractable) reference distribution and the operational $\pi_\text{ref}$ utilized within DPO. Additional implementation specifics and chosen hyperparameters are discussed in Appendix~\ref{app:implementation}.

\section{Theoretical Foundations of DPO} This section delivers a deeper conceptual understanding of the DPO approach, develops supporting theoretical results, and connects the merits of DPO with the known limitations of actor-critic frameworks for RLHF (including algorithms like PPO~\cite{schulman2017proximal}).

\label{sec:theory}

\subsection{A Language Model as an Implicit Reward Model} DPO enables simultaneous reward inference and policy learning through a unified maximum likelihood framework, avoiding the explicit reward estimation and reinforcement learning steps. The loss function in Eq.~\ref{eq:main_eq} matches the Bradley-Terry model structure, employing a reward transformation $r^*(x, y) = \beta \log\frac{\pi^*_\theta(y \mid x)}{\pi_\text{ref}(y \mid x)}$. Optimization proceeds over the parameterized model $\pi_\theta$, analogous to the objective for the reward model shown in Eq.~\ref{eq:reward_model}, modulo a variable substitution. Here, we elaborate the formal underpinnings of this variable transformation, demonstrate that it maintains full generality with respect to recoverable reward functions, and show that it enables retrieval of the optimal policy exactly. To begin, we introduce an equivalence relation among reward functions.

\begin{definition}
Two reward functions $r(x, y)$ and $r'(x, y)$ are said to be equivalent if and only if there exists a function $f$ such that ${r(x, y)-r'(x, y) = f(x)}$.
\end{definition}

This equivalence can be readily verified to satisfy the properties of an equivalence relation, effectively dividing the space of reward functions into equivalence classes. The following lemmas can thus be established:

\begin{lemma}\label{lemma:same_prefrence} Within the context of Plackett-Luce models, and specifically for Bradley-Terry preferences, any two reward functions that belong to the same equivalence class induce identical preference distributions.
\end{lemma}

\begin{lemma}\label{lemma:same_policy}
    Any pair of reward functions in the same equivalence class will yield the same optimal policy in the context of a constrained RL setting.
\end{lemma}

The arguments supporting these lemmas are uncomplicated and are provided in Appendix \ref{app:lemma1}. The first lemma addresses the well-documented under-specification property inherent to the Plackett-Luce family \cite{plackett1975analysis}. Owing to this ambiguity, it is generally necessary to introduce extra identifiability criteria to obtain well-defined MLE estimates for Eq. \ref{eq:reward_model} \cite{bong2022generalized}. According to the second lemma, any reward functions belonging to the same equivalence class necessarily result in identical optimal policies. Consequently, our primary objective is to recover any representative reward function from the correct equivalence class. The following theorem, which we prove in Appendix~\ref{app:thm1}, formalizes this statement:

\begin{theorem}\label{thm:main}
    Subject to mild conditions, every equivalence class of reward functions consistent with the Plackett-Luce (and in particular Bradley-Terry) framework can be expressed via the reparameterization ${r(x, y) = \beta \log \frac{\pi(y\mid x)}{\pi_\text{ref}(y\mid x)}}$, for some model $\pi(y\mid x)$ and a fixed reference distribution $\pi_\text{ref}(y \mid x)$.
\end{theorem}

\begin{sproof}
    Let $r(x, y)$ be an arbitrary reward function, which yields a corresponding optimal policy $\pi_r(y \mid x)$ as specified in Eq. \ref{eq:op_policy}. We now demonstrate that one can construct a member of $r$'s equivalence class using the aforementioned reparameterization. Define the mapping $f$ as
\begin{equation}
    f(r; \pi_\text{ref}, \beta)(x, y) = r(x, y) - \beta\log\sum_{y}\pi_\text{ref}(y\mid x)\exp\left(\frac{1}{\beta}r(x, y)\right)
\end{equation}
Here, $f$ adjusts the reward function by subtracting the logarithm of its partition function, which depends only on $x$. As this normalization only introduces an $x$-dependent term, $f(r; \pi_\text{ref}, \beta)(x, y)$ indeed lies within the equivalence class of $r(x, y)$. Upon substituting $r$ by the right-hand side of Eq.~\ref{eq:main_eq}, which applies to any $r$, we obtain $f(r; \pi_\text{ref}, \beta)(x, y) = \beta \log \frac{\pi_r(y\mid x)}{\pi_\text{ref}(y\mid x)}$. This establishes that $f$ maps $r$ to an equivalent form that fits the proposed reparameterization, demonstrating no loss of generality for our chosen reward model.
\end{sproof}

An alternative interpretation of Theorem~\ref{thm:main} is that it precisely characterizes which reward function, among all those in a given equivalence class, is selected by the DPO reparameterization—namely, the reward function satisfying:

\begin{equation}\label{eq:lag_p}
     \sum_{y}\underbrace{\pi_\text{ref}(y\mid x)\exp\left(\frac{1}{\beta}r(x, y)\right)}_{=\pi(y\mid x)\text{, using Thm.~\ref{thm:main} reparam.}} = 1,
\end{equation}

that is, ensuring $\pi(y\mid x)$ constitutes a proper probability distribution (i.e., is non-negative and sums to one over $y$). Notably, in light of Eq.~\ref{eq:op_policy}, it becomes evident that Eq.~\ref{eq:lag_p} corresponds to the partition function for the optimal policy derived from the reward function $r(x, y)$. A central insight provided by the DPO method is its ability to introduce specific constraints to the otherwise underdetermined Plackett-Luce (and particularly Bradley-Terry) models for preferences, thereby maintaining the expressiveness of the reward function family while also guaranteeing that the optimal policy in Eq. \ref{eq:op_policy} is computationally explicit for every input $x$.

\subsection{Instability of Actor-Critic Algorithms}
This analytical framework is also applicable for understanding instabilities observed with conventional actor-critic methods commonly employed in RLHF, such as PPO. Following the standard RLHF sequence, we focus on the RL fine-tuning phase as detailed in Section \ref{section:prelims}. There is a direct connection here with the control as inference formalism \cite{levine2018reinforcement} when considering the constrained RL formulation in \ref{eq:RL}. Letting the policy be parameterized as $\pi_{\theta}(y\mid x)$, we aim to minimize the divergence $\mathbb{D}_{\text{KL}}[\pi_{\theta}(y|x) \mid \mid \pi^*(y\mid x)]$, where $\pi^*$ denotes the optimal policy obtained in Eq. \ref{eq:optimum_model} via the reward function $r_{\phi}(y, x)$. Upon algebraic manipulation, this minimization gives rise to the following maximization objective:

\begin{equation}\label{eq:AC}
    \max_{\pi_{\theta}}\mathbb{E}_{\pi_{\theta}(y\mid x)}\bigg[\underbrace{r_{\phi}(x, y) -\beta\log\sum_{y}\pi_\text{ref}(y\mid x)\exp\left(\frac{1}{\beta}r_{\phi}(x, y)\right)}_{f(r_{\phi}, \pi_\text{ref}, \beta)} - \underbrace{\beta\log\frac{\pi_{\theta}(y\mid x)}{\pi_\text{ref}(y\mid x)}}_{\text{KL}}\bigg]
\end{equation}

The objective here aligns with that optimized in earlier studies 
\citep{ziegler2020finetuning, stiennon2022learning, bai2022training, ouyang2022training}, which utilize the DPO-equivalent reward pertaining to the class $r_{\phi}$. In this framework, the normalization component within $f(r_{\phi}, \pi_\text{ref}, \beta)$ can be interpreted as representing the soft value function for the reference policy $\pi_\text{ref}$. Although this normalization does not influence the optimality of the final policy, omitting it may result in elevated variance within the policy gradient estimates, consequently destabilizing the learning process. Employing a learned value function to approximate this normalization is possible but introduces its own optimization challenges. As an alternative, past research often normalizes rewards with respect to a human completion baseline, effectively serving as a one-sample Monte-Carlo approximation of the normalization term. Notably, the DPO reparameterization introduces a reward formulation that circumvents the necessity for such baselines entirely.

\section{Experiments}
This section presents a series of empirical analyses investigating DPO's performance in policy optimization directly from human preferences. We begin by exploring a controlled text-generation scenario, focusing on how effectively DPO balances reward maximization and KL-divergence minimization relative to widely used preference-based approaches like PPO. We then examine DPO on larger-scale models and more complex RLHF benchmarks, encompassing tasks such as summarization and dialogue generation. Our results indicate that DPO, with minimal hyperparameter adjustment, matches or outperforms strong alternatives like PPO-based RLHF and strategies that select the best out of $N$ generated samples using a learned reward signal. The experimental methodology is detailed in the following paragraphs, with further procedural information available in Appendix~\ref{app:exp_details}.

\textbf{Tasks.} In our study, we consider three distinct open-domain text generation tasks. Across all settings, the learning algorithms optimize a policy using a dataset of preference examples $\mathcal{D}=\bigl\{x^{(i)}, y_w^{(i)}, y_l^{(i)}\bigr\}_{i=1}^N$. In the case of \textbf{controlled sentiment generation}, the context $x$ is the initial segment of a movie review from the IMDb dataset \cite{maas-EtAl:2011:ACL-HLT2011}. Here, the objective is to generate a completion $y$ exhibiting positive sentiment. To facilitate controlled assessment, we construct pairs of model generations and utilize a pre-trained sentiment classifier to assign preferences such that $p(\text{positive}\mid x,y_w)>p(\text{positive}\mid x,y_l)$. For SFT, we further adapt GPT-2-large on training reviews from IMDb until performance stabilizes (details provided in App~\ref{app:sentiment_details}). In the \textbf{summarization} task, $x$ corresponds to a Reddit forum post and the goal is to generate a summary $y$ highlighting its key points. We follow existing benchmarks, utilizing the Reddit TL;DR dataset \citep{volske-etal-2017-tl} alongside human preference data collected by \citeauthor{stiennon2022learning}. For SFT, a model fine-tuned on summaries authored by humans\footnote{\url{https://huggingface.co/CarperAI/openai_summarize_tldr_sft}} is employed, with the TRLX framework \citep{leandro_von_werra_2023_7790115} applied for RLHF. Note that human preferences were elicited from outputs of a different but comparably trained SFT model. Lastly, for the \textbf{single-turn dialogue} task, $x$ represents a user’s prompt, ranging from scientific queries to requests for personal advice. The task requires a policy to produce a useful and engaging reply $y$; we utilize the Anthropic Helpful and Harmless dialogue dataset \citep{bai2022training}, which comprises 170,000 interactions between a user and an automated system. Each dialogue is paired with two model-generated responses, with human annotation indicating the favored reply. Unlike the previous tasks, no pre-trained SFT model is directly accessible here, so we construct one by fine-tuning a general-purpose language model on the subset of responses judged as preferred.

\textbf{Evaluation.} Our experimental framework employs two primary evaluation strategies. In the context of controlled sentiment generation, where we can access the ground-truth reward function via a sentiment classifier, we assess the performance of each algorithm by mapping out its trade-off curve between the achieved reward and the KL-divergence from the reference policy. This “frontier” serves as a meaningful benchmark for constrained reward maximization. In contrast, such ground truth is unavailable in practical, real-world tasks. To address this, we instead compare each algorithm’s outputs to a baseline using the \textit{win rate}, where GPT-4 acts as a stand-in for human evaluators in assessing summary quality and the helpfulness of responses for the summarization and single-turn dialogue settings, respectively. In summarization, the baseline comprises the test set’s reference summaries, whereas for dialogue, we use the dataset's preferred response as the baseline. Although recent literature indicates that LMs can often serve as more reliable automated judges than standard metrics \citep{Chen2023ExploringTU}, we additionally perform a human subject study to substantiate our reliance on GPT-4 as an evaluator, as discussed in Sec.~\ref{sec:human-judgments}. Our results demonstrate that GPT-4’s decisions align closely with human judgment, often showing inter-rater agreement comparable to or exceeding that among human annotators.

\textbf{Methods.} Beyond DPO, our evaluation spans a variety of established techniques for aligning language model outputs with human preference. For summarization, we experiment with zero-shot prompting using \textbf{GPT-J} \citep{gpt-j}, and for dialogue, we adopt 2-shot prompting with \textbf{Pythia-2.8B} \citep{biderman2023pythia}. We further include the standard \textbf{SFT} model, as well as \textbf{Preferred-FT}, which is constructed by supervised fine-tuning on the preferred output $y_w$: this is derived either from the SFT model in controlled sentiment and summarization, or from a base LM in dialogue tasks. Another approach we examine is \textbf{Unlikelihood}~\citep{welleck2019neural}, which adjusts the model by increasing the likelihood of $y_w$ and explicitly decreasing the likelihood assigned to the disfavored $y_l$, modulated by an optional coefficient $\alpha\in[0,1]$ for the unlikelihood component. Our comparisons also involve \textbf{PPO} \citep{schulman2017proximal}, where a reward model is trained from preference data, as well as an oracle variant, \textbf{PPO-GT}, utilizing the exact reward function available in the controlled sentiment scenario. In our sentiment generation studies, we test two PPO-GT variants: a standard version \cite{leandro_von_werra_2023_7790115} and an enhanced version that normalizes rewards and adjusts hyperparameters for optimality (these refinements are also applied to PPO with learned rewards). As a strong baseline, we incorporate the \textbf{Best of $N$} method, which involves sampling $N$ candidates from the SFT model (or Preferred-FT in the dialogue setting), selecting the best according to a reward function trained from preferences. While this method can yield high performance by separating reward model accuracy from PPO optimization, it is computationally infeasible for moderate values of $N$ because every query at test time requires generating $N$ full outputs.

\subsection{How well can DPO optimize the RLHF objective?}

Standard RLHF algorithms employ a KL-regularized reward maximization framework, which simultaneously encourages the policy to attain high reward and penalizes excessive divergence from a designated reference policy. Consequently, any comparative assessment of such methods must jointly consider both the attained reward and the magnitude of the KL divergence; a modest reward gain does not justify substantially higher KL divergence. In Figure~\ref{fig:frontier-tldr-main}, the reward-KL tradeoff frontiers for various approaches are depicted in the context of the sentiment task. For each method, we perform multiple experiments, adjusting the hyperparameters governing policy conservatism (with target KL values of $\{3,6,9,12\}$ for PPO, $\beta$ set to $\{0.05,0.1,1,5\}$, and $\alpha$ chosen from $\{0.05,0.1,0.5,1\}$ for the unlikelihood approach, and different random seeds for preferred-FT). Across all configurations, this amounts to 22 experimental runs. Following every 100 optimization steps up to convergence, each policy is assessed on a collection of test prompts, calculating both the mean reward under the true reward model and the average sequence-level KL\footnote{That is, the total KL-divergence aggregated across all timesteps.} with respect to the reference policy $\text{KL}\left(\pi\mid \mid \pi_\text{ref}\right)$. Our results demonstrate that DPO yields a significantly more favorable frontier, attaining superior rewards without incurring high KL divergence. This is especially striking for several reasons. Firstly, DPO and PPO target the same underlying objective, yet DPO achieves this with markedly greater efficiency; the reward-to-KL tradeoff for DPO consistently surpasses that of PPO. Secondly, even when PPO operates with access to the true reward signal (PPO-GT), DPO remains able to outperform it on the frontier.

\subsection{Can DPO scale to real preference datasets?}
\label{sec:dpo-real-datasets}
We proceed by assessing DPO’s effectiveness for fine-tuning on real-world tasks, specifically in summarization and single-turn dialogue. In summarization, standard automatic metrics such as ROUGE have been shown to correlate weakly with human judgments~\citep{stiennon2022learning}, and past studies indicate that large language models (LMs) fine-tuned with PPO based on human feedback tend to yield more desirable summaries. To compare the performance of various approaches, we generate completions for the test portion of the TL;DR summarization dataset and determine the mean win rate relative to reference completions in the test set. Sampling for all evaluated methods is performed across temperature values ranging from 0.0 to 1.0; the corresponding win rates are depicted in Figure~\ref{fig:frontier-tldr-main} (right). DPO, PPO, and Preferred-FT are each applied to the same GPT-J SFT model\footnote{\url{https://huggingface.co/CarperAI/openai_summarize_tldr_sft}} for fine-tuning. Results show DPO attains an approximate win rate of 61\% at a sampling temperature of 0.0, surpassing PPO’s best win rate of about 57\% at the same temperature. Furthermore, DPO demonstrates a greater peak win rate than the best-of-$N$ reference baseline. It is worth mentioning that DPO’s $\beta$ hyperparameter was not carefully optimized, suggesting that the current results might underrepresent the approach’s full capability. Additionally, DPO exhibits considerably higher resilience to changes in sampling temperature, in contrast to PPO, whose performance may fall back to that of the unmodified GPT-J model at elevated temperatures. Preferred-FT yields little to no improvement compared to the SFT baseline. Direct human evaluations between DPO and PPO, detailed in Section~\ref{sec:human-judgments}, reveal that DPO outputs sampled at temperature 0.25 are preferred 58\% of the time over PPO samples at temperature 0.

For single-turn dialogue evaluation, we analyze multiple methods on the subset of the Anthropic HH dataset’s test split \citep{bai2022training} that contains a single exchange between human and assistant. When using GPT-4 as an evaluator, the system’s win rate is determined by comparing method outputs against the human-preferred completions in the test set. Since a standardized SFT model does not exist for this benchmark, our pipeline initiates with a pre-trained Pythia-2.8B model. We first apply Preferred-FT to fine-tune a reference model on the selected completions, thereby aligning model outputs with the desired distribution, and subsequently proceed with DPO training. We further benchmark against the strongest outcome from 128 Preferred-FT completions (having observed that performance for the Best of $N$ baseline saturates at $N=128$; refer to Appendix Figure~\ref{fig:best-of-n}) as well as a two-shot prompted variant of the Pythia-2.8B model. Across the most effective sampling temperatures, DPO meets or surpasses the performance of other methods. Additionally, we assess an RLHF model trained via PPO on the Anthropic HH dataset\footnote{\url{https://huggingface.co/reciprocate/ppo_hh_pythia-6B}} from a reputable implementation\footnote{\url{https://github.com/CarperAI/trlx/tree/main/examples/hh}}, but could not identify prompt configurations or sampling temperatures yielding improvements over the unmodified Pythia-2.8B base model. Drawing from observations in the TL;DR task and noting that Best of 128 and PPO both target the same reward function, we treat Best of 128 as an approximate measure of PPO’s efficacy. In sum, DPO emerges as the sole computationally efficient strategy that outperforms the preferred completions within the Anthropic HH dataset, matching or exceeding the results of the resource-intensive Best of 128 baseline. As depicted in Figure~\ref{fig:dialogue-main}, DPO rapidly attains peak performance.

\subsection{Generalization to a new input distribution}

To more rigorously assess the comparative performance of PPO and DPO under shifts in data distribution, we take the policies trained during the Reddit TL;DR summarization experiments and test them on a distinct domain: news articles from the CNN/DailyMail dataset's test split \citep{nallapati-etal-2016-abstractive}. This evaluation uses the optimal sampling temperatures identified on TL;DR (0 and 0.25). The findings are summarized in Table~\ref{tab:ood}. GPT-4 win rates against the ground-truth dataset summaries were measured using the same GPT-4 (C) prompt as with Reddit TL;DR, but substituting ``forum post'' for ``news article''. On this shifted distribution, DPO once again achieves a considerably higher performance than the PPO policy. These results provide preliminary evidence that DPO’s policies not only match but can even surpass PPO’s generalization capabilities, despite the fact that DPO omits the use of additional unlabeled Reddit TL;DR prompts leveraged by PPO.

\subsection{Validating GPT-4 judgments with human judgments}
\label{sec:human-judgments}
A human evaluation study was conducted to assess the reliability of GPT-4’s judgments, leveraging outputs from the TL;DR summarization experiments with two distinct GPT-4 prompts. The \textbf{GPT-4 (S)} (simple) prompt requests a decision regarding which summary better captures the essential information from the original post. The \textbf{GPT-4 (C)} (concise) prompt additionally asks which summary is more succinct; this variation was included because GPT-4’s preferences using the \textbf{GPT-4 (S)} prompt tend toward longer, more repetitive outputs compared to human judgments. Detailed versions of the prompts can be found in Appendix~\ref{app:prompts}. To cover a spectrum of generation quality, three comparisons are carried out: the top-performing (DPO, temp. 0.25), the lowest-performing (PPO, temp. 1.0), and an intermediate model (SFT, temp. 0.25), with each contrasted against greedily-sampled PPO at its best-performing temperature. Results indicate that for both prompts, GPT-4’s preferences are as aligned with human preferences as humans are with each other, indicating that GPT-4 can function as an effective stand-in for human evaluators (multiple human raters were only used for the DPO and PPO-1 pairs, due to constraints on human resources). Among the prompts, \textbf{GPT-4 (C)} was most reflective of human win rates; as such, this prompt was chosen for reporting main results in Section~\ref{sec:dpo-real-datasets}. For further details about the human study, including the rater interface and the roster of volunteers, refer to Appendix~\ref{app:human-study}.

\section{Discussion}
The preference-based learning paradigm offers a robust and adaptable approach for developing proficient, well-aligned language models. In this work, we have presented DPO—a straightforward methodology that enables training language models using preference data, circumventing the complexities of reinforcement learning. Instead of recasting preference learning as a conventional reinforcement learning problem to leverage existing RL methods, DPO establishes a principled correspondence between policy distributions of language models and reward functions, permitting direct optimization toward human preferences using a simple cross-entropy objective. This process entirely eliminates the need for reinforcement learning or reduction in expressive power. Notably, DPO demonstrates comparable or superior performance to established RLHF approaches, such as those utilizing PPO, and achieves this with minimal hyperparameter adjustment, thus substantially lowering the barriers associated with leveraging human preference data for training language models.

\textbf{Limitations \& Future Work.} The findings of this study introduce several open questions for subsequent research. To what extent can the DPO-trained policy generalize to data beyond its training distribution compared to methods that optimize an explicit reward function? Although preliminary evidence points to generalization abilities on par with those achieved by PPO-based techniques, further investigation across diverse scenarios is warranted. For instance, is it possible to leverage unlabeled prompts effectively through self-labeling mechanisms using the DPO policy? Another avenue concerns understanding how reward over-optimization appears in the context of direct preference optimization, and whether the slight performance drop observed in Figure~\ref{fig:dialogue-main}-right is symptomatic of such effects. Furthermore, while this work evaluates models up to 6B parameters, scaling DPO for much larger, cutting-edge models remains an appealing challenge for future research. With respect to evaluation methods, we observe that GPT-4-based win rate assessments are sensitive to the phrasing of prompts; determining optimal strategies for eliciting reliable, high-quality assessments from automated evaluators is a promising area for additional inquiry. Ultimately, DPO presents opportunities extending far beyond language model preference training, suggesting its utility for generative modeling tasks across various data domains.